% Better Poster template (A0 landscape) — Morrison 2019 format adapted for
% Japanese academic posters with bxjsarticle + tcolorbox.
%
% Layout: three vertical zones on A0 landscape (1189 × 841 mm):
%   - Left bar (25%): TL;DR bullets + author + QR
%   - Center billboard (50%): single plain-language main finding, ≥120pt
%   - Right bar (25%): figures, methods, references
%
% Source: Mike Morrison, "How to design a better research poster" (2019).
% Eye-tracking pilot (ScienceUX 2020) showed faster comprehension than
% the traditional dense layout.
\documentclass[a0paper,landscape,25pt]{bxjsarticle}
\usepackage[dvipdfmx]{graphicx}
\usepackage[dvipdfmx]{color,xcolor}
\usepackage{geometry}
\geometry{margin=1.5cm,top=1cm,bottom=1cm}
\usepackage{tcolorbox}
\tcbuselibrary{breakable,skins}
\usepackage{multicol}
\usepackage{enumitem}
\usepackage{anyfontsize}
\usepackage{tikz}
\usepackage{newtxtext,newtxmath}
\renewcommand{\familydefault}{\sfdefault}

% --- Palette: Morrison-style minimalism (3 colors + black/white) ----------
\definecolor{billboardbg}{RGB}{255,255,255}
\definecolor{billboardink}{RGB}{20,20,20}
\definecolor{barbg}{RGB}{0,51,102}
\definecolor{barink}{RGB}{255,255,255}
\definecolor{accent}{RGB}{220,30,30}

\pagestyle{empty}

\begin{document}
\noindent
\begin{minipage}[t]{0.24\linewidth}
    %% LEFT BAR — TL;DR + author + QR
    \begin{tcolorbox}[
        colback=barbg, coltext=barink, boxrule=0pt,
        height=0.97\textheight, valign=top, arc=0pt,
        left=12pt,right=12pt,top=12pt,bottom=12pt
    ]
        {\fontsize{34}{40}\selectfont\bfseries 一目で分かる要点}\\[6pt]
        \rule{\linewidth}{0.5pt}\\[6pt]
        {\fontsize{28}{38}\selectfont
        \begin{itemize}[leftmargin=*, itemsep=4pt]
            \item TL;DR bullet 1 — main mechanism
            \item TL;DR bullet 2 — key result
            \item TL;DR bullet 3 — implication
            \item TL;DR bullet 4 — limitation
        \end{itemize}}

        \vspace{20pt}
        \rule{\linewidth}{0.5pt}\\[10pt]
        {\fontsize{24}{30}\selectfont\bfseries 著者}\\[4pt]
        {\fontsize{22}{28}\selectfont 近畿大学 理工学部 電気電子工学科\\菅原賢悟}

        \vspace{20pt}
        \rule{\linewidth}{0.5pt}\\[10pt]
        {\fontsize{24}{30}\selectfont\bfseries フルペーパー}\\[6pt]
        \begin{center}
        \fbox{\parbox[c][45mm][c]{45mm}{\centering [QR code\\here]}}\\[4pt]
        {\fontsize{18}{22}\selectfont Scan for full paper}
        \end{center}
    \end{tcolorbox}
\end{minipage}\hfill
\begin{minipage}[t]{0.48\linewidth}
    %% CENTER BILLBOARD — single plain-language finding, ≥120pt
    \begin{tcolorbox}[
        colback=billboardbg, coltext=billboardink, boxrule=0pt,
        height=0.97\textheight, valign=center, arc=0pt,
        left=20pt,right=20pt,top=20pt,bottom=20pt
    ]
        \begin{center}
        {\fontsize{120}{140}\selectfont\bfseries
        \textcolor{accent}{Main finding}\\
        in plain language}\\[40pt]
        {\fontsize{50}{60}\selectfont
        Subtitle: 1 sentence elaboration that\\
        a non-specialist can understand}
        \end{center}
    \end{tcolorbox}
\end{minipage}\hfill
\begin{minipage}[t]{0.24\linewidth}
    %% RIGHT BAR — figure / method / refs
    \begin{tcolorbox}[
        colback=billboardbg, coltext=billboardink, boxrule=1pt,
        colframe=barbg,
        height=0.97\textheight, valign=top, arc=0pt,
        left=12pt,right=12pt,top=12pt,bottom=12pt
    ]
        {\fontsize{30}{36}\selectfont\bfseries タイトル / Title}\\[4pt]
        {\fontsize{22}{28}\selectfont サブタイトル}\\[12pt]
        \rule{\linewidth}{0.5pt}\\[10pt]

        {\fontsize{26}{32}\selectfont\bfseries 手法}\\[4pt]
        {\fontsize{22}{28}\selectfont 簡潔な手法説明 (3-5 行)。専門用語は最小限に。}

        \vspace{12pt}
        \rule{\linewidth}{0.5pt}\\[10pt]

        {\fontsize{26}{32}\selectfont\bfseries 結果図}\\[6pt]
        \begin{center}
        \fbox{\parbox[c][80mm][c]{80mm}{\centering [Key figure]}}\\[4pt]
        {\fontsize{20}{24}\selectfont 図1: 結果の図 (軸+単位+凡例)}
        \end{center}

        \vspace{12pt}
        \rule{\linewidth}{0.5pt}\\[10pt]

        {\fontsize{26}{32}\selectfont\bfseries 参考文献}\\[4pt]
        {\fontsize{18}{22}\selectfont
        \begin{enumerate}[label={[\arabic*]}, leftmargin=*, itemsep=0pt]
            \item Ref 1
            \item Ref 2
        \end{enumerate}}
    \end{tcolorbox}
\end{minipage}

\end{document}
